\section{Tartışma}\label{sec:discussion}

\paragraph{Hangi durumda hangi çözücü} PSB, hedeflerin sırasını onları robotlara atamadan önce
sabitlemekte, ardından optimal bölme rotaları eşik sınaması başına $O(nk)$ sürede dengelemektedir;
kurulum maliyeti $n$ ile neredeyse doğrusal olarak artmakta ve ilk çözüm saniyenin küçük bir
kesrinde elde edilmektedir. Literatürde öne çıkan
metasezgiseller ve LKH-3 çok daha geniş bir çözüm uzayını taramakta ve yeterli süre verildiğinde
daha iyi çözümler bulmaktadır; özellikle MILS, küçük örneklerde birkaç on saniye, orta büyüklükteki
örneklerde ise yaklaşık iki dakika sonra PSB+LNS'yi geçmiştir. Saniyeler içinde tamamlanması
gereken planlamada, binlerce hedef içeren örneklerde ve daha uzun bir aramaya başlangıç
çözümü sağlamak için PSB daha uygun bir seçenektir; birkaç dakikalık sürenin bulunduğu orta büyüklükteki örneklerin
çevrimdışı planlanmasında ise MILS daha uygundur.

\paragraph{Bahçede önemli olan} Hasat benzetiminde görev süresinin büyük kısmını toplama süresi oluşturmaktadır;
bu nedenle rotalama farkları görev süresine yüzde birkaçtan yaklaşık yüzde ona kadar uzanan
farklar olarak yansımaktadır. En büyük etkiler, çözücünün arama gücünden bağımsız üç karardan
kaynaklanmaktadır: toplama geçişi sırasında ortaya çıkan meyve sayılarıyla yeniden planlama,
bahçenin kendi ekseninde planlama (PSB için yaklaşık 12\%) ve planlama süresinin kasa
değişimleri arasındaki aralığın altında tutulması. PSB, plan başına milisaniyeler düzeyindeki süresiyle son
koşulu sağlamaktadır; 1\,s süre sınırıyla çalışan OR-Tools PSB'ye 2--3\% yaklaşmış, 10\,s sınırla ise
1\,ha'lık bahçelerde PSB ile benzer ortalama sonuç vermiştir. Uç devralmalı merkezi olmayan $k$-d
bölüntüsü, merkezi yeniden planlamanın mümkün olmadığı durumlarda pratik bir alternatiftir
ve devralma, bu yaklaşımın kaybının önemli bir kısmını telafi etmektedir.

\paragraph{Sınırlılıklar} (i) Hasat sonuçları benzetimden elde edilmiştir; saha denemesi
yapılmamıştır ve kablolu platform parametreleri sektör basınındaki tanımlamalara dayanan
varsayımlardır. Duyarlılık analizi yalnızca Tablo~\ref{tab:params}'de aralık verilen
parametreleri kapsar; durma/dönüş süresi gibi diğer varsayımlar ayrıca sınanmamıştır.
(ii) Algılama, görüntüler üzerinde bir dedektör çalıştırılarak değil, istatistiksel olarak
(yakalama oranı, kesinlik, görünür pay) modellenmiştir. (iii) Platformlar arasındaki etkileşimler
modellenmemiştir (çarpışma ve kablo kısıtı yoktur, kasa istasyonu kapasitesi sınırsızdır) ve
hareket, bahçe ekseninde $L_1$ uzaklıklarıyla yaklaşık olarak hesaplanmıştır; bölgede yaygın olan,
farklı yönelimli birden çok küçük parselden oluşan işletmeler ayrıca modellenmemiştir.
(iv) Literatürde öne çıkan kodlar tek bir tohum ve varsayılan parametrelerle çalıştırılmış,
bunlardan ikisi yalnızca yürütülebilir dosya olarak kullanılabilmiştir; MA en büyük örnekte
çalıştırılamamıştır. Ortak \emph{best} referansı farklı sayıda tohum içeren yöntemlerin en
iyisinden hesaplandığından sapmalar tanımlayıcıdır; 20 eşzamanlı koşunun 12 fiziksel
çekirdekteki bellek ve zaman paylaşımı etkileri ayrıca izole edilmemiştir.
(v) Alt sınır ara $k$ değerleri için zayıftır; bu nedenle bahçelerdeki mutlak kalite yalnızca
yaklaşık olarak bilinmektedir. (vi) Yöntemler eşit hasat miktarında durdurulmamıştır; bu
nedenle süre, olgun meyve payı ve verimlilik birlikte değerlendirilmelidir. Kümelenmiş
güven aralıkları yalnızca 8--10 bahçe--geçiş kümesine dayanır. (vii) Görselleştirme,
benzetimin yeniden oynatımıdır; bir uçuş testi değildir.
